\section{Grammatical Framework}\label{app_sec:grammatical_framework}

Grammatical Framework (GF)~ is a programming language for multilingual grammar applications.
It is a special-purpose language for grammars, like YACC, Bison, Happy, BNFC, but not restricted to programming languages.
It is a functional programming language, like Haskell, Lisp, OCaml, SML, Scheme, but specialized to grammar writing.
For example, the following is a GF grammar to generate simple English sentences about food:

\begin{lstlisting}[language=haskell,caption={The abstract syntax of the Food grammar.},captionpos=b,label={lst:gf food grammar}, basicstyle=\small\ttfamily]
abstract Food = {
  flags startcat = Comment ;
  cat
    Comment ; Item ; Kind ; Quality ;
  fun
    Pred : Item -> Quality -> Comment ;
    This, That : Kind -> Item ;
    Mod : Quality -> Kind -> Kind ;
    Wine, Cheese, Fish : Kind ;
    Very : Quality -> Quality ;
    Fresh, Warm, Italian,
}
Expensive, Delicious, Boring : Quality ;
\end{lstlisting}


\subsection{Abstract Syntax vs. Concrete Syntax}\label{subsec:abstract-syntax-vs-concrete-syntax}

Constrained decoding works by pruning the set of allowed tokens-id at each decoding step.
This requires that both the grammar and the completion engine work at the token-id level.
Since the tokenization scheme of LMs are usually different from each other, the grammar becomes dependent on the tokenization scheme of the LM.
%
This brings two challenges: (1) the grammar needs to be redefined for each LM (tokenization scheme), and (2) the debugging of the grammar is difficult because the grammar is defined at the token-id level.
%
We propose to decouple the grammar from the tokenization scheme of the LM by defining an abstract grammar and a couple of concrete grammars.
%
The abstract grammar is defined at the text level, which is human-readable and independent of the tokenization scheme of the LM.
%
The concrete grammars are defined at the token-id level, which are dependent on the tokenization scheme of the LM and work with the completion engine.
%
Once an abstract grammar is defined, the concrete grammars can be automatically translated from the abstract grammar.
%
This is similar to the process of compiling a high-level programming language to a low-level assembly language.
%
The separation of the abstract grammar and the concrete grammar is implemented in GF.
%
For example, the predication rule

\begin{lstlisting}[language=haskell,label={lst:gf abs},caption={BNF Notation},captionpos=b, basicstyle=\small\ttfamily]
  Pred. Comment ::= Item "is" Quality
\end{lstlisting}

now becomes two rules:

\begin{lstlisting}[language=haskell,label={lst:gf crt},caption={Grammatical Framework Notation},captionpos=b, basicstyle=\small\ttfamily]
  fun Pred : Item -> Quality -> Comment ;
  lin Pred item quality = item ++
            "is" ++ quality ;
\end{lstlisting}

All that matters in a \textbf{linearization rule} is that it defines a string as a function of the variables that it depends on.
%
A GF grammar consists of two parts: \textbf{abstract syntax} and \textbf{concrete syntax}.Below is the concrete syntax of the aforementioned Food grammar:

\begin{lstlisting}[language=haskell,label={lst:gf crt food grammar},caption={The concrete syntax of the Food grammar.},captionpos=b, basicstyle=\small\ttfamily]
  concrete FoodEng of Food = {
  lincat
    Comment, Item, Kind, Quality = Str ;
  lin
    Pred item quality = item ++ "is"
        ++ quality ;
    This kind = "this" ++ kind ;
    That kind = "that" ++ kind ;
    Mod quality kind = quality ++ kind ;
    Wine = "wine" ;
    Cheese = "cheese" ;
    Fish = "fish" ;
    Very quality = "very" ++ quality ;
    Fresh = "fresh" ;
    Warm = "warm" ;
    Italian = "Italian" ;
    Expensive = "expensive" ;
    Delicious = "delicious" ;
    Boring = "boring" ;
}
\end{lstlisting}

And also the concrete syntax of the Food grammar in Italian:

\begin{lstlisting}[language=haskell,label={lst:gf crt food grammar italian},caption={The concrete syntax of the Food grammar in Italian.},captionpos=b, basicstyle=\small\ttfamily]
concrete FoodIta of Food = {
    lincat
      Comment, Item, Kind, Quality = Str ;
    lin
    Pred item quality = item ++ "e"
        ++ quality ;
    This kind = "questo" ++ kind ;
    That kind = "quel" ++ kind ;
    Mod quality kind = kind ++ quality ;
    Wine = "vino" ;
    Cheese = "formaggio" ;
    Fish = "pesce" ;
    Very quality = "molto" ++ quality ;
    Fresh = "fresco" ;
    Warm = "caldo" ;
    Italian = "italiano" ;
    Expensive = "caro" ;
    Delicious = "delizioso" ;
    Boring = "noioso" ;
}
\end{lstlisting}

\subsection{Multilingual Grammars}\label{subsec:multilingual-grammars}

A \textbf{multilingual grammar} is a system with \textbf{one} abstract syntax and \textbf{any number} of concrete syntaxes.
%
While GF was originally designed for machine translation, it happens to be well suited to handle multi-tokenizations. Different large language models have different tokenizers. For the same sentence, their representations in token id space are different. This phenomenon is very similar to the multilingual grammar scenario, where the same abstract syntax tree can be linearized into different languages as shown in in the Food grammar.

\subsection{Expressivity}\label{subsec:expressivity}

Grammatical Framework (GF)'s incremental parser, supports PMCFG (Parallel Multiple Context-Free Grammars).
PMCFG lies between mildly context-sensitive and fully context-sensitive grammars~\cite{SEKI1991191}.


